%% kapa: exam=1 round=31 year=2020 period=2 form=A subject=accounting subject_ko=회계학 seq=40 number=80
%% source: Q-Net 감정평가사 기출문제 (cst00309) artlSeq=5207219 file=extracted/1008/5207219/A형/A형/2020년도 제31회 감정평가사 1차 시험_2교시_A형.pdf
%% conversion: kordoc 4.24.51 -> tools/gen1.py (auto); review=proofread-sonnet; flags=none
\begin{question}{80}
레저용 요트를 전문적으로 생산ㆍ판매하고 있는 (주)감평은 매년 해당 요트의 주요 부품인 자동제어센서 2,000단위를 자가제조하고 있으며, 관련 원가자료는 다음과 같다.
\par{\small\centering\begin{tabular}{|c|c|c|}\hline
구 분 & 총원가 & 단위당원가 \\ \hline
직접재료원가 & ￦700,000 & ￦350 \\ \hline
직접노무원가 & 500,000 & 250 \\ \hline
변동제조간접원가 & 300,000 & 150 \\ \hline
고정제조간접원가 & 800,000 & 400 \\ \hline
합계 & ￦2,300,000 & ￦1,150 \\ \hline
\end{tabular}\par}
(주)감평은 최근 외부업체로부터 자동제어센서 2,000단위 전량을 단위당 ￦900에 공급하겠다는 제안을 받았다. (주)감평이 동 제안을 수락할 경우, 기존설비를 임대하여 연간 ￦200,000의 수익을 창출할 수 있으며, 고정제조간접원가의 20\%를 회피할 수 있다. (주)감평이 외부업체로부터 해당 부품을 공급받을 경우, 연간 영업이익에 미치는 영향은?
\begin{choices}
  \item ￦0
  \item ￦60,000 감소
  \item ￦60,000 증가
  \item ￦140,000 감소
  \item ￦140,000 증가
\end{choices}
\end{question}
\answer{3}
